%%%%%%%%%%%%%%%%%%%%%%%%%%%%%%%%%%%%%%%%%%%%%%%%%%%%%%%%%%%%%%%%%%
%                           Document style
%%%%%%%%%%%%%%%%%%%%%%%%%%%%%%%%%%%%%%%%%%%%%%%%%%%%%%%%%%%%%%%%%%
\documentclass[12pt]{report}

% Importing report style and the first pages
%%%%%%%%%%%%%%%%%%%%%%%%%%%%%%%%%%%%%%%%%%%%%%%%%%%%%%%%%%%%%%%%%%
%                       Main document style
%%%%%%%%%%%%%%%%%%%%%%%%%%%%%%%%%%%%%%%%%%%%%%%%%%%%%%%%%%%%%%%%%%

% Importing user packages (librarys for different functions and commands)
\usepackage[utf8]{inputenc}
\usepackage{subfiles}               % add subfiles
\usepackage[T1]{fontenc}            % text font
\usepackage{lmodern}                % font
\usepackage{pdfpages, pdflscape}    % add pdf's
\usepackage{gensymb, textcomp}      % add symbols
\usepackage{lastpage}               % get the last page number
\usepackage{graphicx}               % add images
\usepackage{setspace}               % define size of space (infinite)
\usepackage{xcolor}                 % colors
\usepackage{comment}                % comments
\usepackage{fancyhdr}               % to get fancy headers
\usepackage[background=black, arrow=red, text=white]{callouts}
\usepackage{wrapfig}                % positioning figures
\usepackage{subcaption}             % position mulitple figures
\usepackage[toc,page]{appendix}     % add appendix

% hungarian language
\usepackage[utf8]{inputenc}
\usepackage{t1enc}
\def\magyarOptions{defaults=hu-min}
\usepackage[magyar]{babel}

% page settings
\setlength{\marginparwidth}{2cm}    % margins
\usepackage[a4paper,left=2.5cm,right=2.5cm,top=2.8cm,bottom=4cm]{geometry}
\setlength\parindent{0pt}
\setlength{\parskip}{1em}

\setlength{\headheight}{37pt} % header hight

\cfoot\thepage % page number as footer

% creating base header
\pagestyle{fancy}
\rhead{\includegraphics[height=0.7cm]{./images/simonyi-logo.png}}                   % itt a magasságán lehet finomhangolni

%define colors for links
\definecolor{ceruleanblue}{rgb}{0.16, 0.32, 0.75}
\definecolor{coolblack}{rgb}{0.0, 0.18, 0.39}


% clickable table of contents
\usepackage{hyperref}
\hypersetup{linktoc = all, colorlinks=true,
    linkcolor=coolblack, filecolor=magenta,
    urlcolor=ceruleanblue, citecolor=ceruleanblue}


%%%%%%%%%%%%%%%%%%%%%%%%%%%%%%%%%%%%%%%%%%%%%%%%%%%%%%%%%%%%%%%%%%
%                   Command for the first page
%%%%%%%%%%%%%%%%%%%%%%%%%%%%%%%%%%%%%%%%%%%%%%%%%%%%%%%%%%%%%%%%%%

\newcommand*\FirstPage{
\thispagestyle{empty}
% placing a pdf as background and adding text on top
\includepdf[pages=1,picturecommand*={
    % selecting placement, font size and text type
    \put(215,630){
        \fontsize{28}{28}\selectfont{
            \textbf{SZERVEZETI ÉS}
        }
    }
    \put(140,590){
        \fontsize{28}{28}\selectfont{
            \textbf{MŰKÖDÉSI SZABÁLYZAT}
        }
    }
    \put(208,540){                                              % Itt az X koordinátát lehet finomhangolni
        \fontsize{28}{28}\selectfont{
            \Large{\textbf{Simonyi Károly Szakkollégium}}
        }
    }
    \put(140,50){
        \fontsize{28}{28}\selectfont{
            \Large{\textbf{Budapest, 2024. október 1.}}
        }
    }
}]{./images/simonyi-frontpage-bg.pdf}}


% section header
\lhead{Szervezeti és Működési Szabályzat - Simonyi Károly Szakkollégium}
\linespread{1.25}

%%%%%%%%%%%%%%%%%%%%%%%%%%%%%%%%%%%%%%%%%%%%%%%%%%%%%%%%%%%%%%%%%%
%                       New Commands
%%%%%%%%%%%%%%%%%%%%%%%%%%%%%%%%%%%%%%%%%%%%%%%%%%%%%%%%%%%%%%%%%%

%creating paragraphs --> \newpara
\newcommand{\newpara}{\vskip 2mm}

%add pdf's --> \addpdf
\newcommand*\addpdf[1]{
\includepdf[pages=-]{#1}}

%add pdf's in landscape --> \addpdflandscape
\newcommand*\addpdflscape[1]{
\begin{landscape}
\newgeometry{margin=0cm}
\includepdf[pages=-,angle=90]{#1}
\end{landscape}}

%add sections with 4 numbers --> \paragraph
\usepackage{titlesec}
\setcounter{secnumdepth}{4}
\titleformat{\paragraph}
    {\normalfont\normalsize\bfseries}{\theparagraph}{1em}{}
\titlespacing*{\paragraph}
    {0pt}{3.25ex plus 1ex minus .2ex}{1.5ex plus .2ex}

\titleformat{\section}
    {\normalfont\Large\bfseries}{\thesection.\S}{1em}{} 
\renewcommand{\thesection}{\arabic{section}}
\renewcommand{\thesubsection}{\Roman{subsection}}


% sections without numbers, but still in table of contents
%\setcounter{secnumdepth}{0} % remove "%" to include

%=================================================================
%                           Start Document
%=================================================================
\begin{document}

% adding the first two pages
\FirstPage
\pagebreak

\section{A Simonyi Károly Szakkollégium Szervezeti és Működési Szabályzata}

\begin{enumerate}
  \item Jelen szabályzat a Budapesti Műszaki és Gazdaságtudományi Egyetem Villamosmérnöki és Informatikai Karon (továbbiakban: BME VIK) belül 
    működő Simonyi Károly Szakkollégium (továbbiakban Szakkollégium) Szervezeti és Működési Szabályzata (továbbiakban: SZMSZ). A Szakkollégium 
    befogadó szerve a BME VIK Hallgatói Képviselete.
  \item A Simonyi Károly Szakkollégium a Schönherz Zoltán Szakkollégium és a Kollégiumi Önkormányzat Szakmai Reszortjának jogutódja.
  \item A Szakkollégium a Szakkollégiumi Mozgalom által megalkotott Szakkollégiumi Chartát magáénak tekinti.
  \item A Szakkollégium az Oktatási Hivatal által nyilvántartott intézmény.
\end{enumerate}

\section{A Szakkollégium adatai}

\begin{enumerate}
  \item \textbf{A Szakkollégium neve}: Simonyi Károly Szakkollégium
  \item \textbf{Idegen nyelvű (angol) neve}: Simonyi Karoly College for Advanced Studies
  \item \textbf{Rövidítése}: Simonyi
  \item \textbf{Elérhetősége}: info@simonyi.bme.hu
  \item \textbf{Címe}: Schönherz Kollégium, 1117 Budapest, Irinyi József utca 42.
  \item \textbf{Alapítás éve}: 2003.
\end{enumerate}

% Section border ==========================================================

\section{A Szakkollégium célja}

\begin{enumerate}
  \item A Szakkollégium célja, hogy elsősorban a Szakkollégium tagjai és egyetemi hallgatók számára lehetőséget biztosítson az egyetemi képzést 
    kiegészítő, magasabb szintű ismeretek elsajátítására tanfolyamok és előadások szervezésével. Továbbá lehetőségeihez mérten infrastruktúrával, 
    mentorálással és konzultációk szervezésével támogatja a szakkollégisták önálló szakmai tevékenységeit.
  \item A Szakkollégium céljának tekinti, hogy a tagjai által elért eredményeket minél szélesebb körben publikálja, a magyar felsőoktatás más 
    egységeivel és szakemberekkel megismertesse a Szakkollégiumban folyó munkát.
  \item Továbbá célja, hogy az általa összefogott szakmai köröknek és csoportoknak szervezett lehetőséget biztosítson a színvonalas működéshez; 
    összefogja, koordinálja és felügyelje azokat, illetve érdekeiket képviselje a lehető legszélesebb körben.
  \item A Szakkollégium mindenkori céljának tekinti a szakkollégiumi hagyományok ápolását és fejlesztését.
  \item A Szakkollégium céljának tekinti más hazai és külföldi szakkollégiumokkal, szakmai szervezetekkel való kapcsolatfelvételt és együttműködést.
\end{enumerate}

% Section border ==========================================================

\section{A Szakkollégium felépítése}

\begin{enumerate}
  \item A Szakkollégiumot az általa befogadott, a Szakkollégium elveit és jelen szabályzatot elfogadó szakmai körök (továbbiakban: szakterületek) és 
    csoportok alkotják. A szakterületek aktuális listáját jelen szabályzat \ref{annex_currents} melléklete tartalmazza. A csoportokra vonatkozó szabályokat jelen 
    szabályzat \ref{annex_groups} melléklete tartalmazza.
  \item A Szakkollégium a Schönherzes Villamosmérnökök és Informatikusok Egyesületének (továbbiakban: SVIE) Alapszabálya által definiált reszort szerepét 
    tölti be, és ennek megfelelően képviseli a Szakkollégiumot alkotó szakterületeket, koordinálja munkájukat.
  \item A Szakkollégium megfelel a Budapesti Műszaki és Gazdaságtudományi Egyetem Hallgatói Önkormányzat (továbbiakban: BME HÖK) Alapszabályának és a 
    szakkollégiumokat érintő mellékleteinek.
  \item A Szakkollégium tagja a Műegyetemi Szakkollégiumok Közösségének (továbbiakban: MŰSZAK), amelynek működését a BME HÖK Alapszabályának 9. számú 
    melléklete (A Műegyetemi Szakkollégiumok Közösségének Szervezeti és Működési Szabályzata) szabályozza. 
  \item A Szakkollégium fő szervei az elnökség, Szervezeti Bizottság (továbbiakban: SZB) és a vezetőség.
  \item A Szakkollégium döntéshozó szerve a vezetőség, amely a szakterületek vezetőiből, az SZB tagjaiból, az elnökségből, valamint a vezetőség 
    \ref{sec_leadership}. § \ref{enum_leadership_voting}. alapján megválasztott csoportvezetőkből áll.
  \item A Szakkollégium munkáját a Szakkollégium elnöke koordinálja.
  \item A Szakkollégium működését segíti és a szabályos működését ellenőrzi a Felügyelő Bizottság (továbbiakban: FEB).
  \item A Szakkollégiumhoz tartozó szakterületek és csoportok belső szabályzataikat maguk alkotják, de nem állhatnak jelen szabályzattal ellentmondásban.
\end{enumerate}

% Section border ==========================================================

\section{A Szakkollégiumi tagság} \label{sec_membership}

\begin{enumerate}
  \item A szakkollégiumi tevékenységekben a tagjelöltek, a rendes tagok, az archivált tagok, valamint az öregtagok vesznek részt.
  \item A szakkollégiumi tagság önkéntes.
  \item A Szakkollégium nem szed tagdíjat a tagjaitól.
\end{enumerate}

\subsection{Tagjelöltek}

\begin{enumerate}
  \item A Szakkollégiumba tagjelöltnek jelentkezhet bármely felsőoktatásban résztvevő hallgató.
  \item Jelentkezni a Szakkollégiumot alkotó valamelyik szakterületnél lehet az adott szakterület saját szabályzata alapján.
  \item A jelentkezést az adott szakterület bírálja el a saját szabályzatában rögzítettek szerint.
  \item Több szakterülethez való jelentkezés esetén ezek elbírálása egymástól függetlenül történik.
  \item A tagjelölti viszony megszűnik, ha az adott szakterület a saját szabályzatában rögzített módon megszünteti, a tagjelölt saját kérésére, vagy fegyelmi eljárás eredményeképp.
\end{enumerate}

\subsection{Rendes tagok} \label{subsec_member}

\begin{enumerate}
  \item A Szakkollégium rendes tagjai elfogadják és magukra nézve kötelezőnek ismerik el a Szakkollégium SZMSZ-ét és belső szabályzatait.
  \item A szakkollégiumi rendes tagok a tagjelöltek sorából kerülhetnek ki. Rendes taggá csak az válhat, aki a Szakkollégiumot alkotó szakterületek valamelyikének tagjelöltje.
  \item Rendes taggá válni a szakkollégiumi felvételi eljárás keretében lehetséges. A felvételi eljárás mindenkori leírását jelen szabályzat \ref{annex_admission} melléklete tartalmazza.
  \item A Szakkollégium minden félévben általános felvételi eljárást hirdet. A felvételi mindenkori általános elveit jelen szabályzat \ref{annex_admission} melléklete tartalmazza. 
    A mindenkori felvételi kritériumokat az adott félévben a vezetőség határozza meg.
  \item A rendes tagok egyes szakterületekben végzett munkája minden félév végén értékelésre kerül az adott szakterületi vezető által.
  
  \item A szakkollégiumi rendes tagság megszűnik:
  \begin{enumerate}
    \item ha a szakkollégiumi rendes tag minden szakkollégiumi szakterületben való tagsága megszűnik,
    \item ha a szakkollégiumi rendes tag nem felel meg a \ref{annex_quality} melléklet \ref{subsec_quality_regular}. \ref{enum_quality_regular} pontjában leírtaknak,
    \item fegyelmi eljárás eredményeképp,
    \item a hallgatói jogviszony megszűnésével, kivéve, ha öregtaggá nyilvánítását előterjesztették,
    \item kilépési szándék jelzésével.
  \end{enumerate}

  \item A rendes tagok jogai: \label{enum_member_rights}
  \begin{enumerate}
    \item a Szakkollégium  szolgáltatásait, helyiségeit és eszközeit a vezetőség jóváhagyásával igénybe veheti,
    \item a Szakkollégium belső tanfolyamait, előadásait látogathatja,
    \item a Szakkollégium által szervezett egyéb szakmai, illetve közösségi rendezvényeken részt vehet.
  \end{enumerate}

  \item A rendes tagok kötelezettségei jelen szabályzat \ref{annex_quality} mellékletében olvasható.
\end{enumerate}

\subsection{Archivált tagok}

\begin{enumerate}
  \item Archivált taggá nyilvánítható az a rendes tag, aki szakmai munkát végzett a Szakkollégiumban.
  \item Az archivált taggá nyilvánítást a vezetőség végzi.
  \item Az archivált tagok a Szakkollégium rendes tagjaival azonos jogokat élveznek az \ref{sec_membership}. § \ref{subsec_member}. \ref{enum_member_rights}. pontja alapján.
  \item A szakkollégiumi archív tagság megszűnhet fegyelmi eljárás eredményeképp.
\end{enumerate}

\subsection{Öregtagok}

\begin{enumerate}
  \item Öregtaggá nyilvánítható az a rendes tag, aki kiemelkedő szakmai munkát végzett a Szakkollégiumban.
  \item Az öregtaggá nyilvánítást a vezetőség végzi.
  \item Az öregtagok a Szakkollégium rendes tagjaival azonos jogokat élveznek az \ref{sec_membership}. § \ref{subsec_member}. \ref{enum_member_rights}. pontja alapján.
  \item A szakkollégiumi öregtagság megszűnhet fegyelmi eljárás eredményeképp.
\end{enumerate}

\subsection{Fegyelmi eljárás}

\begin{enumerate}
  \item Fegyelmi eljárást rendes tag javaslatára a FEB indíthat.
  \item A FEB feladata kijelölni a fegyelmi bizottságot, melynek összetételéről a vezetőség dönt többségi szavazással.
  \item A fegyelmi eljárás eredményéről a fegyelmi bizottság dönt.
  \item A fegyelmi eljárás eredményével szemben a fegyelmi bizottság fellebbezési lehetőséget biztosít. A fegyelmi eljárásban érintett 
    ezen szándékát - a közléstől, ennek hiányában a tudomására jutásától számított 15 napon belül - a fegyelmi@simonyi.bme.hu címre küldött e-mailben jelezheti.
  \item Amennyiben a fegyelmi eljárásban érintett fellebbezéssel él, a fegyelmi bizottság köteles a fegyelmi eljárás eredményét újratárgyalni.
\end{enumerate}

% Section border ==========================================================

\section{A Szakkollégium tevékenységei}

\begin{enumerate}
  \item A szakkollégiumi szakmai tevékenységek közé tartozik az elméleti és gyakorlati oktatás, projekt és felkérés alapú szakmai munka.
  \item A szakkollégium egyéb tevékenységei közé tartozik az infrastruktúra üzemeltetése, a szervezetfejlesztés és a rendezvényszervezés.
  \item A szakkollégiumi tevékenységekben a tagjelöltek, a rendes tagok, az archivált tagok, valamint az öregtagok vesznek részt.
  \item A szakterületen folyó szakmai tevékenység minőségéért és annak ellenőrzéséért a megfelelő szakterületi vezető a felelős.
  \item A szakkollégiumban folyó egyéb tevékenységek minőségéért és annak ellenőrzéséért a szakkollégium elnöke a felelős.
  \item A további szabályokat jelen szabályzat \ref{annex_quality} melléklete határozza meg.
\end{enumerate}

% Section border ==========================================================

\section{Szakterületi vezetők} \label{sec_leaders}

\begin{enumerate}
  \item Szakmai kör szakterületi vezetője csak BME VIK-es nappali tagozatos hallgatói jogviszonnyal rendelkező szakkollégiumi rendes tag lehet.
  \item A szakmai körök szakterületi vezetőjének választását és visszahívását a szakterület Szervezeti és Működési Szabályzata határozza meg.
  \item A szakterületi vezető feladata a szakterület szabályzatában meghatározottakon felül:
  \begin{enumerate}
    \item értékelni a projektek résztvevőinek munkáját a projektvezetők beszámolója alapján,
    \item ellenőrizni a projektek szakmai színvonalát,
    \item képviselni a szakterület érdekeit a vezetőségi gyűléseken,
    \item képviselni a vezetőség döntéseit a szakterülete felé,
    \item a vezetőség bármely tagjának kérésére beszámolni a hozzá tartozó feladatkörökben végzett munkájáról,
    \item a Szakkollégium stratégiai terveinek kidolgozásában és megvalósításában szerepet vállalni.
  \end{enumerate}
\end{enumerate}

% Section border ==========================================================

\section{A Szakkollégium elnöke} \label{sec_president}

\begin{enumerate}
\item A Szakkollégium elnöke csak BME VIK-es hallgatói jogviszonnyal rendelkező szakkollégiumi rendes tag lehet.
\item A Szakkollégium elnökének mandátuma egy naptári évre szól, amely megválasztását követő szeptember 1-jétől a következő év augusztus 31-éig tart.
\item Az elnök \label{enum_pres_tasks}

\begin{enumerate} 
  \item képviseli a Szakkollégiumot a Szakkollégiumot érintő fórumokon,
  \item tartja a kapcsolatot a Szakkollégium partnereivel,
  \item ellenőrzi a szakkollégiumi tevékenységek szabályzat szerinti működését,
  \item figyeli a pályázati lehetőségeket, és gondoskodik a pályázatok tartalmi és formai követelményeknek megfelelő elkészítéséről, beadásáról, utánkövetéséről és beszámolásáról,
  \item szervezi a Szakkollégium közösségi életét és újoncok toborzását,
  \item segíti utódja betanulási folyamatát, utódja számára biztosítja az akadálymentes posztátadást,
  \item \ref{sec_leadership}. § \ref{enum_leadership_meeting}. pontja alapján összehívja a vezetőségi gyűléseket, a napirendi pontokat számon tartja, a gyűlést levezeti,
  \item \ref{sec_presidency}. § \ref{enum_presidency_meeting}. pontja alapján összehívja az elnökségi üléseket, a napirendi pontokat számon tartja, az ülést levezeti,
  \item a Szakkollégium felvételi és minőségbiztosítási folyamatait indítványozza, felügyeli és levezeti,
  \item a vezetőségi tagok munkájával kapcsolatosan beszámoltatást kezdeményez,
  \item felelős a Szakkollégium gazdálkodásáért,
  \item felelős a Szakkollégium közös infrastrukturális eszközeinek rendeltetésszerű használatáért és a Szakkollégium közös tulajdonának megóvásáért,
  \item felelős a csoportokban folyó tevékenység minőségéért és annak ellenőrzéséért.  
\end{enumerate}

\item A Szakkollégium elnöke a \ref{enum_pres_tasks}. pontban felsorolt feladataival más szakkollégiumi rendes tago(ka)t vagy öregtag(ok)at is megbízhat.
\end{enumerate}

% Section border ==========================================================

\section{Elnökség} \label{sec_presidency}

\begin{enumerate}
  \item A Szakkollégium elnöke legalább két alelnököt kell, hogy válasszon maga mellé.
  \item A Szakkollégium elnökségét a Szakkollégium elnöke és alelnökei alkotják.
  \item Alelnöki pozíciót kizárólag rendes tag tölthet be.
  \item Az elnökségi pályázatot a FEB írja ki olyan határidőkkel, hogy a választás és annak sikertelensége esetén a pótpályáztatás 
    és a pótválasztás megtörténhessen a jelenlegi elnökség mandátumának lejárta előtt. A pályázási időszak legalább két hét. 
    A pályázási időszak vége és az elnökségválasztás között legalább két hétnek kell eltelnie.
  \item A Szakkollégium elnökségének személyét az érvényes pályázatot vagy pótpályázatot leadók közül választják meg kétharmados többséggel a választáson, 
    szükség esetén a pótválasztáson. A választáson szavazásra jogosultak a vezetőség \ref{sec_leadership}. § \ref{enum_leadership_voting}. pont alapján szavazati joggal 
    rendelkező vezetőségi tagok az elnökség kivételével. A szavazásra jogosult akadályoztatás esetén kiválaszthat egy, a szakterülethez vagy csoporthoz tartozó rendes tagot, 
    aki a személyében eljár. Amennyiben az összeférhetetlenség lehetősége fennáll, a szavazásra jogosult köteles egy, a szakterületéhez vagy csoportjához 
    tartozó rendes tagot kiválasztani, aki a személyében eljár.
  \item A vezetőség bármely tagja indítványozhatja a FEB bevonásával az elnökség visszahívását. Amennyiben a vezetőség szavazati jogú tagjainak legalább 
    kétharmada támogatja a visszahívást, az elnökséget le kell mondatni, ebben az esetben a FEB köteles eljárni a \ref{enum_presidency_resign}. pont alapján.
  \item Amennyiben a mindenkori elnökség lemond, visszahívják vagy nem teljesíti az elnökségi pozíciójának betöltéséhez szükséges feltételeket, 
    és már rendelkezik a szakkollégium újonnan megválasztott elnökséggel, őt kell beiktatni ügyvezetőként mandátuma megkezdéséig. 
    Ellenkező esetben a további szabályokról a \ref{enum_presidency_temp}. pont rendelkezik. \label{enum_presidency_resign}
  \item Amennyiben a szakkollégium nem rendelkezik újonnan megválasztott elnökséggel és az elnökség lemond, visszahívják vagy nem teljesíti az elnökségi 
    pozíciójának betöltéséhez szükséges feltételeket, illetve sikertelen a választás és pótválasztás, ebben az esetben a FEB ügyvezetőt jelöl ki a 
    vezetőség aktuális tagjai közül, aki átmenetileg a következő elnökség megválasztásáig ellátja az elnökségi feladatokat. A lemondástól, visszahívástól, 
    vagy a legutóbbi sikertelen választástól számított 5 munkanapon belül a FEB megkezdi az időközi elnökségválasztást. \label{enum_presidency_temp}
  \item Az időközi elnökségválasztás a FEB által kiírt időközi elnökségi pályázattal kezdődik. A pályázatot olyan határidőkkel kell meghirdetni, hogy az 
    időközi elnökségválasztás 3 hét alatt lezajlhasson.
  \item Az időközi elnökségválasztás során megválasztott elnökség mandátuma megválasztásával kezdődik és megválasztása utáni augusztus 31-éig tart.
  \item Az időközi elnökségválasztás procedúráját annyiszor kell megismételni, amíg nem kerül megválasztásra elnökség. Két egymást követő sikertelen 
    időközi elnökségválasztás után a vezetőség kell kiállítson maguk közül egy elnökséget, akinek mandátuma kiállításával kezdődik és kiállítása utáni augusztus 31-éig tart.
  \item Az alelnökök az elnök akadályoztatása esetén minden tekintetben helyettesítik őt.
  \item Az elnökség a vezetőségben együttesen egy szavazattal rendelkezik. \label{enum_presidency_voting}
  \item Az elnökség tagjai a \ref{sec_president}. § \ref{enum_pres_tasks}. pontjában meghatározott feladatokat maguk között osztják fel.
  \item Az elnökség tagjainak kötelessége beszámolni munkájukról az elnökségi üléseken és a vezetőségi gyűléseken, a vezetőségi gyűléseken és az elnökségi üléseken részt venni.
  \item Az elnökség a döntéseit az ülésein hozza. Az elnökségi üléseket az elnök hívja össze, szervezi és vezeti le. \label{enum_presidency_meeting}
\end{enumerate}

% Section border ==========================================================

\section{Szervezeti Bizottság} \label{sec_org_committee}

\begin{enumerate}
  \item Szervezeti bizottsági pozíciót kizárólag rendes tag tölthet be.
  \item Az elnök pályázatot ír ki az elnökség terheinek csökkentése céljából a Szakkollégium igazgatása szempontjából fontos feladatok és feladatkörök ellátására.
  \item A határidőig beérkezett, a pályázati kiírásnak formailag is megfelelő pályázatokat a vezetőség bírálja el és dönt róluk többségi szavazással.
  \item A megszavazott személyek alkotják a Szervezeti Bizottságot (továbbiakban: SZB).
  \item Az SZB tagjai a vezetőség nem szavazati jogú tagjai.
  \item Az SZB tagok az elnökségi üléseken tanácskozási joggal rendelkeznek.
  \item Az SZB tagjait az elnök az elnökségi ülések látogatására kötelezheti.
  \item A Szakkollégiumnak mindenkor rendelkeznie kell egy gazdasági felelőssel, ami SZB pozíciónak minősül. Ezen kívül az elnökség belátása szerint SZB pozíciókat hozhat létre 
    vagy szüntethet meg a vezetőség jóváhagyásával.
  \item Egy személy több SZB pozíciót is betölthet. 
  \item Az SZB tagjai az elnökségi üléseken számolnak be munkájukról az elnökség számára.
  \item Az SZB tagjainak munkáját az elnökség felügyeli.
  \item Az SZB tagjait az vezetőség visszahívhatja.  
\end{enumerate}

% Section border ==========================================================

\section{Vezetőség} \label{sec_leadership}

\begin{enumerate}
  \item A vezetőség feladata a Szakkollégium szervezeti kérdéseiben való állásfoglalás, valamint ezzel kapcsolatos döntések meghozatala.
  \item A vezetőség a döntéseit a vezetőségi gyűléseken hozza.
  \item A vezetőségi gyűléseken kötelező részt vennie a szakterületek vezetőinek, az elnökségnek valamint a \ref{enum_leadership_voting}. alapján szavazati joggal rendelkező csoportvezetőknek.
  \item A vezetőségi gyűlések szabadon látogathatók bármely szakkollégiumi rendes tag, archivált tag és öregtag számára, amennyiben a vezetőség ezzel ellentétes,
    az adott témára vonatkozó döntést nem hoz. Személyeket érintő döntések vendégek jelenléte nélkül tárgyalandók.
  \item Az elnök köteles a vezetőségi gyűléseket a szorgalmi időszakban heti rendszerességgel összehívni, amennyiben külső tényezők által nem akadályoztatott. \label{enum_leadership_meeting}
  \item A szakterületi vezetők legalább 50\%-ának írásos kérésére a vezetőségi gyűlést össze kell hívni.
  \item A vezetőség szavazati jogú tagjai az elnökség a szabályzat \ref{sec_presidency}. § \ref{enum_presidency_voting}. pontja alapján, a szakterületek vezetői valamint amennyiben egy csoport 
    által végzett tevékenység megköveteli, a vezetőség szavazással dönthet arról, hogy a csoportvezető szavazati joggal rendelkező vezetőségi taggá váljon. A vezetőség szavazati joggal 
    rendelkező tagjai egyetlen szavazattal rendelkeznek. A szakterület- vagy csoport vezetője akadályoztatás esetén kiválaszthat egy, a szakterülethez vagy csoporthoz 
    tartozó rendes tagot, aki a személyében eljár, a kiválasztott személyről a vezetőség szavazással dönt. \label{enum_leadership_voting}
  \item A vezetőség határozatképes, ha a szavazati joggal rendelkező tagjainak legalább fele jelen van.
  \item A vezetőség határozatait - az elnök személyére vonatkozó döntések kivételével - egyszerű többséggel hozza.
  \item A vezetőség feladata jóváhagyni a Szakkollégium beadandó és kiírandó pályázatait, továbbá meghatározni a Szakkollégium partnereivel való együttműködés feltételeit.
\end{enumerate}

% Section border ==========================================================

\section{Felügyelő Bizottság} \label{sec_supervisory_board}

\begin{enumerate}
  \item A FEB tagjainak mandátuma a következő FEB beiktatásáig szól. A szakkollégium két évente választja meg a FEB tagjait a \ref{enum_feb_election}. pont alapján.
  \item A FEB három tagból áll.
  \item A FEB tagjait a vezetőség kéri fel és szavazza meg, törekedve arra, hogy a szavazást megtartsák az adott naptári év első vezetőségi gyűlésén. FEB tagságra a 
    vezetőség tagjai jelölhetnek öregtagokat. A három legtöbb szavazatot kapó személy kerül felkérésre, aki a szavazatok több mint felét megkapta. \label{enum_feb_election}
  \item A FEB legalább egy tagjának lemondása esetén a vezetőség időközi FEB választást kell végrehajtson a lemondást követő két héten belül.
  \item A FEB feladatai:
  \begin{enumerate}
    \item Iránymutatást adni a Szakkollégium elnökségének és vezetőségének.
    \item A Szakkollégiumban folyó szakmai és egyéb tevékenységek minőségének és szabályosságának ellenőrzése.
    \item A Szakkollégium elnökségének félévenkénti beszámoltatása. Ez történhet írásban vagy szóban. Indokolt esetben rendkívüli beszámolót kérhetnek.
    \item A FEB tagjai jóváhagyják az elnökség beszámolóját. A jóváhagyáshoz a három FEB tagból minimum kettőnek a szavazata szükséges.
    \item Amennyiben az elnökségi beszámolót a FEB nem hagyja jóvá, kezdeményezheti rendkívüli vezetőségi gyűlés összehívását melynek témája az elnökség visszahívása. 
      Rendkívüli vezetőségi gyűlésen a vezetőség kétharmados többségi szavazata esetén az elnökség visszahívásra kerül, a továbbiakban a \ref{sec_presidency}. § \ref{enum_presidency_resign}. pontja alapján köteles eljárjon a FEB.
    \item Jóváhagyja a megválasztott elnökség személyét. A jóváhagyáshoz a három FEB tagból minimum kettőnek a szavazata szükséges. Amennyiben a FEB nem hagyja jóvá az 
      elnökség személyét, ebben az esetben a \ref{sec_presidency}. § \ref{enum_presidency_resign}. pontja alapján köteles eljárjon a FEB.
    \item A FEB egy adott évben két különböző leendő elnökséget utasíthat vissza, a harmadikat köteles jóváhagyni.
    \item Jóváhagyja az SZMSZ módosításokat, amennyiben azok valóban a szervezet fejlődését szolgálják. A jóváhagyáshoz a három FEB tagból minimum kettőnek a szavazata 
      szükséges. Ha a FEB nem hagyja jóvá az SZMSZ módosítását, úgy a vezetőségnek újra kell azt tárgyalnia. Az újratárgyaláson a FEB is részt vesz, a vezetőség által megszavazott 
      új változatot a FEB-nek jóvá kell hagynia.
    \item A fegyelmi eljárás lefolytatása.
  \end{enumerate}
\end{enumerate}

% Section border ==========================================================

\section{A Szakkollégium gazdálkodása}

\begin{enumerate}
  \item A Szakkollégium a vonatkozó felsőbb, egyetemi, kari és kollégiumi szabályzatokban rögzített elvek szerint és módon gazdálkodik.
  \item A Szakkollégium egyéb költségeit pályázatok útján is biztosíthatja.
\end{enumerate}

% Section border ==========================================================

\section{Záró rendelkezések}

\begin{enumerate}
  \item Jelen szabályzatban nem érintett kérdésekben az egyetemi, a kari és a kollégiumi szabályzatok, valamint más ide vonatkozó jogszabályok az irányadóak.
  \item Indokolt esetekben jelen szabályzattól való eltérésről a vezetőség kétharmados szavazással dönt.
  \item A szabályzat mellékletei a következők:
  \begin{enumerate}
    \item \ref{annex_currents} melléklet: A Szakkollégiumot alkotó körök és csoportok aktuális listája
    \item \ref{annex_admission} melléklet: Szakkollégiumi felvételi
    \item \ref{annex_quality} melléklet: Minőségbiztosítási kézikönyv
    \item \ref{annex_groups} melléklet: A Szakkollégium csoportjainak működése
  \end{enumerate}
  \item Jelen szabályzat értelmezésére a FEB jogosult.
  \item Jelen szabályzatot a vezetőség [dátum] hagyta jóvá.
  \item Jelen szabályzatot a BME VIK Hallgatói Képviselete [dátum] hagyta jóvá.
  \item A mellékletek módosítása a Szakkollégium Vezetőségének hatásköre. Módosítási javaslatot a Vezetőség bármely tagja előterjeszthet.
  \item Jelen szabályzat [dátum] lép hatályba. Ezzel egy időben hatályát veszti a szabályzatban érintett kérdésekkel foglalkozó minden korábbi szabályzat. Ezen szabályzatokat 180 napon belül frissíteni kell.
\end{enumerate}

%=================================================================
%                           Appendix
%=================================================================

\appendix
\renewcommand\thechapter{\arabic{chapter}.}
\renewcommand\appendixname{melléklet}

\chapter{A Szakkollégiumot alkotó körök és csoportok aktuális listája} \label{annex_currents}

\begin{enumerate}
  \item \textbf{Hang-, fény- és rendezvénytechnika}: AC Studio \& Live
  \item \textbf{Videotechnika}: Budavári Schönherz Stúdió
  \item \textbf{Rádiótechnika}: HA5KFU
  \item \textbf{Web- és mobilalkalmazás fejlesztés}: Kir-Dev
  \item \textbf{Robotika és MI}: LEGO Kör
  \item \textbf{Digitális látványtervezés}: schdesign
  \item \textbf{Elektronika}: Schönherz Elektronikai Műhely
  \item \textbf{Rendezvény- és műtermi fényképezés}: SPOT
\end{enumerate}

% Chapter border ==========================================================

\chapter{Szakkollégiumi felvételi} \label{annex_admission}

A szakkollégiumi felvételi eljárás célja a tagjelöltek rendes tagként felvétele jelen szabályzatban meghatározott elvek alapján.

\section{A Szakkollégiumba felvételt nyerhet}

\begin{enumerate}
  \item A BME VIK bármely hallgatója (vagy a vezetőség külön engedélyével hallgatói jogviszonnyal rendelkező, felsőoktatási intézményben tanuló hallgató), 
    aki a Szakkollégiumhoz tartozó szakterületek valamelyikével tagjelölti viszonyban áll és a felvételi eljárás során a felvételi kritériumoknak eleget tesz.
\end{enumerate}

\section{Felvételi Bizottság}

\begin{enumerate}
  \item A vezetőség minden felvételi eljárás során megválaszt egy Felvételi Bizottságot (továbbiakban: FB), aminek tagjai lehetnek: nevelőtanárok, a FEB tagjai, 
    a vezetőség tagjai, illetve a vezetőség által javasolt egyéb szakkollégiumi rendes és öreg tagok.
  \item Az FB-t minden felvételi eljárás során újraválasztják. Ennek a választásnak legkésőbb a felvételi időszak kezdő időpontja előtt egy héttel tartott 
    vezetőségi gyűlésen meg kell történnie.
  \item Az FB feladata:
  \begin{enumerate}
    \item a szóbeli beszélgetésben való részvétel,
    \item a felvételi eredményének kiértékelése a meghatározott elvek alapján.
    \item Az FB-nek nem feladata
    \item a felvételi elveinek meghatározása, ami a vezetőség jogköre,
    \item a felvételi kritériumok meghatározása, ami szintén a vezetőség jogköre.
  \end{enumerate}
  \item Az FB tagjainak legalább három különböző szakterületről kell kikerülniük. Az esetlegesen felmerülő stratégiai vagy elvi kérdésekben kizárólag a vezetőség dönthet.
\end{enumerate}

\section{A felvételi menete}

\begin{enumerate}
  \item A Szakkollégium minden félévben felvételit hirdet. A felvételi időszak ütemezését az elnökség határozza meg, úgy, hogy az legkésőbb egy héttel a felvételi 
    időszak megkezdése előtt kihirdethető legyen. Az FB feladata a felvételi folyamat lebonyolítása és adminisztrálása. A felvételi eredményét az elnök és a 
    szakterületi vezetők hirdetik ki a mindenkori Félévnyitó és Felvételi Vacsorán.
\end{enumerate}

\section{A felvételi alapelvei}

\begin{enumerate}
  \item \textbf{Szakmai alkalmasság}: A szakterületek a jelentkezők szakmai alkalmasságáról szakmai teszt vagy beadandó feladat formájában bizonyosodhatnak meg.
  \item \textbf{Szakterületi tagjelöltként végzett szakmai munka}: A szakterületi vezetők kötelesek arról tájékoztatni az FB-t, hogy a felvételizők a Szakkollégiumban 
    töltött idejük alatt a szükséges mennyiségű és minőségű szakmai munkát elvégezték-e. A fentebb említett kritériumokról pontosabb utasításokat jelen szabályzat 3. számú melléklete tartalmazza.
  \item \textbf{Célkitűzés és gondolkodásmód}: Meg kell győződni arról, hogy a jelentkező tisztában van-e a Szakkollégium céljaival és a szakkollégiumi gondolkodásmóddal, 
    valamint arról, hogy mit vár a Szakkollégiumtól illetve milyen formában tervezi segíteni a Szakkollégium működését és fejlődését.
  \item A jelentkező ilyen jellegű álláspontját irányított szóbeli beszélgetésben vagy motivációs levél jellegű írásbeli felméréssel lehet megállapítani. 
    Ha a felmérés szóban zajlik, a feltett kérdés csoportokat abban az esetben is előre ki kell dolgoznia az FB-nek, és azt a vezetőségnek jóvá kell hagynia. 
    Javasolt, hogy a kérdések térjenek ki eddigi tanulmányi eredményekre is, hiszen ezek fontos részét képezik a felvételizőkkel szembeni elvárásoknak.
\end{enumerate}

% Chapter border ==========================================================

\chapter{Minőségbiztosítási kézikönyv} \label{annex_quality}

\section{A kézikönyv célja}

\begin{enumerate}
  \item A kézikönyv célja, hogy lefektesse a Szakkollégium minőségbiztosítási alapelveit.
\end{enumerate}

\section{Az alkalmazott módszer}

\begin{enumerate}
  \item A minőségbiztosítás alapja a szakkollégiumban végezhető tevékenységek részletes leírása.
  \item A Szakkollégiumnak alapvetően két dolgot kell a minőségbiztosítás irányítása alá helyeznie: az egyéni fejlődést és a lehetőséget a szakmai fejlődésre,
    vagyis a Szakkollégiumban végzett szakmai tevékenységek tematikus rendszerének az ápolását.
\end{enumerate}

\section{Szakkollégiumban végezhető tevékenységek fajtái} \label{sec_quality_activities}

\begin{enumerate}
  \item Egyéni tanulás, munka
  \begin{enumerate}
    \item Alapozó kurzusaink mindenki számára nyitottak, de az igazi fejlődés a szakterülethez csatlakozva érhető el, tehát a szakkollégiumi taggá válás nem a végső cél, 
      hanem a kezdet. A tagoktól elvárjuk, hogy szakmai projektmunka mellett tovább fejlesszék magukat.
    \item Fontosnak tartjuk a szakmai fejlődésen kívül minden olyan képesség megszerzését is, amely mind a karrierben, mind a civil életben való helytállást biztosíthatja.
  \end{enumerate}

  \item Projektmunka
  \begin{enumerate}
    \item A szakkollégiumban végezhető munkák javarészt projektek. A projektek sokfélék lehetnek:
    \begin{enumerate}
      \item Meghatározott időre lezáruló projekt: sajátossága, hogy előre meghatározott időn belül be kell fejezni, illetve beszámolni az eredményekről. Pl.: rendezvény, 
        közvetítés, stúdiós felvételek, weboldal design készítése stb.
      \item Fejlesztési projekt: a fejlesztés specifikálásával kezdődik és nem meghatározható időn belül ér véget. Erről félévente kell beszámolni a szakterületi vezetőjének. 
        Pl.: tanfolyamrobot fejlesztése, stúdiófelújítás.
      \item Tervezési projekt: rövid/közép/hosszú távú, stratégiai tervek meghatározása. Erre a szakterületi vezető kérheti fel a tagot. Erről való beszámolás eseti 
        megegyezés alapján történik.
      \end{enumerate}
  \end{enumerate}

  \item Oktatási tevékenység
  \begin{enumerate}
    \item A Szakkollégium fő tevékenysége az oktatás, így nagy figyelmet kell fordítani annak minőségére. Az oktatási tevékenységnél kiemelten fontos az oktatott téma 
      részletes kidolgozása, előadása, és az oktatásban résztvevők értékelése.
  \end{enumerate}

  \item Kutatás
  \begin{enumerate}
    \item A kutatás olyan projekt, amelynek célja egy probléma vizsgálata, aminek eddig nem ismert a megoldása. A kutatási munka tervezéssel, a kutatás lépéseinek 
      specifikálásával kezdődik. A kutatásról félévente kell beszámolni a szakterületi vezetőnek, aki indokolt esetben dönthet a felfüggesztés mellett.
  \end{enumerate}

  \item Üzemeltetői feladatok
  \begin{enumerate}
    \item Az üzemeltetői feladat folyamatos munkavégzés, mely a Schönherz Kollégiumban működő rendszerek működésének felügyeletét, karbantartását és a felmerülő hibák javítását 
      foglalja magában. A felmerülő problémáról és annak megoldásáról értesíteni kell a szakterületi vezetőt, és félévente be kell számolni az elvégzett munkáról, illetve a fejlesztési lehetőségekről.
  \end{enumerate}
\end{enumerate}

\section{Szakkollégiumi tagok fejlődési lehetőségei és kötelezettségei}

\subsection{Tagjelölt}

\begin{enumerate}
  \item Szakterületi feladatai:
  \begin{enumerate}
    \item Fejlődés, adott szakterület tanfolyamain való részvétel
    \item Felkészülés a taggá válásra
    \item Projektekben való részvétel
    \item Általa választott szakterület tagjainak megismerése, közösségi életébe való bekapcsolódása
  \end{enumerate}

  \item Szakkollégiumi feladatai:
  \begin{enumerate}
    \item Simonyi Károly Szakkollégium felépítésének megismerése
    \item A Szakkollégium történetének és céljainak megismerése
    \item A Szakkollégium közösségi és szakmai rendezvényein való részvétel
  \end{enumerate}
\end{enumerate}

\subsection{Rendes tag} \label{subsec_quality_regular}

\begin{enumerate}
  \item Szakterületi feladatai:
  \begin{enumerate}
    \item Adott szakmai kör vagy csoport közösségi és szakmai munkájában való aktív részvétel
    \item Fejlődés (évente minimum egy tanfolyamon, tréningen vagy projektben való részvétel)
    \item Beszámolás félévente
  \end{enumerate}

  \item Szakkollégiumi feladatai: \label{enum_quality_regular}

  \begin{enumerate}
    \item Szakmai kritérium: Ezt a kritériumot azon tagok teljesítik, akik a dokumentum \ref{sec_quality_activities}. § pontjaiban felsorolt tevékenységi körökben,
      azaz a szakkollégiumi munkában aktívan részt vettek az adott félévben. A kritérium teljesítésének megállapítása a tag által választott szakterületi- vagy 
      csoportvezetőjének hatásköre és feladata. Az eredmények jelentése a vezetőség felé minden félévben a vizsgaidőszak végéig esedékes. 
    \item Közösségi kritérium: A közösségi feltételnek azok felelnek meg, akik az adott félévben minimum két, a szakkollégium által szervezett rendezvényen 
      részt vettek. A kritérium teljesítésének megállapítását az elnökség végzi.
    \item Tanulmányi kritérium: A Szakkollégium tanulmányi kritériuma a 3.0-s kreditindex megléte minden félév végén. A tanulmányi kritérium ellenőrzése az 
      elnökség feladata.
  \end{enumerate}

  \item A minőségbiztosításon való megfeleléshez a fenti három kritériumból a szakmai kritériumot kötelezően teljesítenie kell a rendes tagnak, 
    a másik kettőből egynek szükséges teljesülnie. Ez alól felmentést kaphat a rendes tag szakkollégiumi pályafutása során egy alkalommal szakterületi 
    vezetőjének javaslatára a vezetőség jóváhagyásával. A tagság megszűnésével elvesznek a Szakkollégiumi jogok és kötelességek. A tagság megszűnése után 
    ismételten taggá csak akkor válhat, ha újrafelvételizik.
  \item A tagság megszűnéséről a Szakkollégium elnöke, vagy az általa megbízott szakterületi vezető emailben értesíti a tagot. A címzett az email 
    elküldése utáni egy hétben kérelmezheti kizárásának felülvizsgálatát a Szakkollégium elnökének írt email formájában.   
\end{enumerate}

\section{Szakkollégiumi rendes és öregtagként vállalható posztok}

\subsection{Nevelőtanár}
\begin{enumerate}
  \item A nevelőtanárrá válás feltételei:
  \begin{enumerate}
    \item az adott szakterület és a vezetőség által is elismert munkát végzett,
    \item rendelkezik legalább egy BSc diplomával,
    \item személyének támogatásáról a vezetőség szavazással dönt.
  \end{enumerate}
  \item A nevelőtanár feladatai:
  \begin{enumerate}
    \item a folyamatos tanácsadás biztosítása a Szakkollégiumban,
    \item a vezetőség tagjainak mentorálása és motiválása,
    \item konfliktusok feloldásának ösztönzése.
  \end{enumerate}
\end{enumerate}

\section{Fegyelmi ülés}

\begin{enumerate}
  \item Amennyiben valamely kizárt tag vagy archivált tag a meghatározott időn belül jelezte, hogy felül szeretné vizsgáltatni az ügyét, az elnöknek a fellebbezések határideje 
    után három héten belül meg kell hirdetnie egy fegyelmi ülést - amely lehet egy vezetőségi ülés vendégek nélkül tárgyalandó része -, amelyen a fellebbezők szakterületi vezetőjük 
    jelenlétében megvédhetik a tagságukat.
  \item Minden ügy kimeneteléről a vezetőség dönt. Egy-egy ügy tárgyalásának két kimenetele lehet:
  \begin{enumerate}
    \item Tagság megújítása: A kizárt szakkollégista visszaszerzi tagságát, ezzel pedig minden jogát és kötelezettségét. Ehhez a vezetőség kétharmados jóváhagyása szükséges. 
    \item Elutasítás: Amennyiben a szakkollégista tagsága nem kerül megújításra, nem élhet további fellebbezéssel.
  \end{enumerate}
\end{enumerate}

\section{Tagság szüneteltetése}

\begin{enumerate}
  \item Egy tag szüneteltetheti tagsági viszonyát maximum két félévre a szakterületi vezetőjének írásban benyújtott kérelemmel, melyet a vezetőség hagyhatja jóvá. 
    Ez alatt az idő alatt nem kell teljesítenie a tagság feltételeit, illetve nem élhet a tagság nyújtotta előnyökkel. A kért idő letelte után automatikusan aktívvá válik, 
    és teljesíteni kell a tagsági feltételeket.
\end{enumerate}

\section{Simonyi Plakett}

\begin{enumerate}
  \item Simonyi Plakettben részesülhet az a szakkollégiumi öregtag, aki szakmai körének, illetve csoportjának kiemelten megbecsült tagja, szakkollégiumi pályafutása arról ad számot, 
    hogy a szakkollégium tagjai elé példaként állítható.
  \item A Plakettre való jelöléshez jelölési időszakot indítványozhat a vezetőség bármely tagja, ekkor az elnökség a jelölési időszakot megnyitja, a végén a jelöltek Plakettben való 
    részesítését személyenként a vezetőség többségi szavazással elbírálja.
  \item A Simonyi Plakettet a Félévnyitó és Felvételi Vacsorán ünnepélyes keretek között adják át a jelölők. A nyertesek a Plakett mellé egy szakmai munkát igazoló iratot kapnak, 
    melyet a szakterületi vezető, illetve az elnök ellenjegyez.
\end{enumerate}

% Chapter border ==========================================================

\chapter{A Szakkollégium csoportjainak működése} \label{annex_groups}

\section{Csoport működése}

\begin{enumerate}
  \item A csoportnak joga van részt venni a szakkollégiumot érintő pályázatokban.
\end{enumerate}

\section{Csoport alapítása}

\begin{enumerate}
  \item A csoportnak rendelkeznie kell öt leendő taggal.

  \item A csoportnak rendelkeznie kell egy működésére vonatkozó szabályzattal, aminek része kell, hogy legyen:
  \begin{enumerate}
    \item a csoport célja,
    \item csoport tevékenysége,
    \item csoportvezető megválasztásának módja,
    \item taggá válás feltételei,
    \item tagság megszűnésének folyamata.
  \end{enumerate}

  \item A csoport megalapításához kérelmet kell benyújtani az elnökség felé, a kérelemnek tartalmaznia kell:
  \begin{enumerate}
    \item a csoport nevét,
    \item csoportvezető nevét,
    \item alapítótagok listáját,
    \item a csoport működésére vonatkozó szabályzatot.
  \end{enumerate}
  \item A szabályzat elfogadásáról és a csoport megalakulásáról a vezetőség szavazással dönt.
\end{enumerate}

\section{Csoport mgszűnése}

\begin{enumerate}
  \item A csoport megszűnését a csoport indítványozhatja.
  \item A csoport megszűnéséről a vezetőség szavazással dönt.
  \item Amennyiben nem rendelkezik csoportvezetővel legalább három hónapig, abban az esetben meg kell szüntetni.
  \item Amennyiben a csoport célját és működését nem látja el, abban az esetben az elnökségnek indítványoznia kell a csoport megszüntetését.
\end{enumerate}

\section{Csoport felépítése}

\begin{enumerate}
  \item Csoportvezető
  \begin{enumerate}
    \item felelős a csoport által végzett tevékenységért,
    \item feladata a csoport szabályzatának betartása és betartatása,
    \item a vezetőség vagy az elnökség kérésére köteles beszámolni a csoport tevékenységéről.
  \end{enumerate}

  \item Tag
  \begin{enumerate}
    \item Kizárólag rendes szakkollégiumi tag lehet,
    \item elfogadja a csoport szabályzatát.
  \end{enumerate}
\end{enumerate}

\end{document}
